\section{A closed-form global game of stablecoin redemption}\label{sec:model}

This section formalises the stress depeg piece of the holding leg as a global game in the closed-form tradition of \citet{klee_etal_2026_fragility}. Uniform signal noise makes the redeeming mass linear in the fundamental, so the equilibrium threshold solves a single one-dimensional condition and the expected depeg loss, the stress probability, and the reserve comparative static are analytical. The model prices the one piece of the circuit no observable conversion cost reveals; the counterparty and baseline pieces are read from data in Section~\ref{sec:risk_factors}. Proofs are in Appendix~\ref{app:proofs}.

\subsection{Primitives and information}\label{sub:model_primitives}

There is a continuum of homogeneous holders of unit mass. The latent issuer fundamental is $\theta$, drawn from a normal prior $\theta \sim \mathcal{N}(1, \sigma_\theta^2)$ concentrated at par: the issuer is healthy almost surely, and a run occurs only in the lower tail. Each holder $i$ observes a private signal with \emph{uniform} noise,
\begin{equation}\label{eq:signal}
s_i = \theta + \sigma \varepsilon_i, \qquad \varepsilon_i \sim \mathcal{U}[-1, 1],
\end{equation}
so the posterior $\theta \mid s_i$ is uniform on $[s_i - \sigma, s_i + \sigma]$. The issuer redeems at par up to liquid reserves $R$; the secondary market has linear price impact $\psi > 0$. The per-redemption transaction cost is $\phi$, the convenience yield of holding the coin is $\omega$, and the counterparty term is $\xi h$, read from market proxies in Section~\ref{sub:counterparty}.

\subsection{Redemption and the secondary price}\label{sub:model_secondary}

The equilibrium is in threshold strategies, redeem if and only if $s_i < s^*$. Under uniform noise the mass that redeems at fundamental $\theta$ is \emph{linear} on the transition interval,
\begin{equation}\label{eq:D}
D(\theta; s^*) = \tfrac{1}{2} + \frac{s^* - \theta}{2\sigma}, \qquad \theta \in [s^*-\sigma,\, s^*+\sigma],
\end{equation}
equal to $1$ below the interval and $0$ above. Redemption runs through two channels: each holder first redeems at par up to reserves $R$, and unmet demand $U(D) = \max\{D - R,\, 0\}$ spills into the secondary market, which clears at
\begin{equation}\label{eq:psec}
p^{sec}(D) = \max\!\left\{\underline{p},\; 1 - \psi\, U(D)\right\}.
\end{equation}
The price falls linearly in unmet redemption relative to depth $1/\psi$, with $\underline p$ a floor. The boundary at which the secondary first opens, $D = R$, is reached at the fundamental
\begin{equation}\label{eq:theta_s}
\theta_s = s^* - \sigma\,(2R - 1),
\end{equation}
below which the coin trades below par. Linearity is the modelling choice that buys the closed form; the convex generalisation is numerically close and is discussed in Appendix~\ref{app:proofs}.

\subsection{Payoffs and the one-dimensional threshold}\label{sub:model_equilibrium}

Conditional on $\theta$ and the aggregate $D$, the payoff to redeeming is par on the share met from reserves and the depressed secondary price on the rest,
\begin{equation}\label{eq:VR}
V_R(D) = \min\!\left\{1, \tfrac{R}{D}\right\} + \left(1 - \min\!\left\{1, \tfrac{R}{D}\right\}\right) p^{sec}(D) - \phi,
\end{equation}
and the payoff to holding is
\begin{equation}\label{eq:VH}
V_H(\theta) = \theta + \omega - \xi h.
\end{equation}
The counterparty term $\xi h$ enters the value of \emph{holding}, not the realised price: netted from the price it would be common to both actions and cancel in the indifference condition, leaving the threshold unmoved; in $V_H$ it lowers the value of holding and feeds the coordination channel. The marginal holder at $s_i = s^*$ is indifferent in expectation over the uniform posterior $\theta \mid s^* \sim \mathcal{U}[s^*-\sigma, s^*+\sigma]$,
\begin{equation}\label{eq:F}
F(s^*) := \mathbb{E}_{\theta \mid s^*}\!\left[V_R(D(\theta; s^*)) - V_H(\theta)\right] = 0.
\end{equation}
This is a single equation in the single unknown $s^*$. The posterior is local and does not involve the fundamental prior, so the threshold is a Laplacian object pinned by the payoff primitives $(R, \psi, \omega, \phi, \xi h, \sigma)$; the prior $\mathcal{N}(1,\sigma_\theta^2)$ enters only the unconditional frequency of each $\theta$, and hence the stress probability and the expected loss below. Equation~\eqref{eq:F} carries a single logarithmic term from the $R/D$ prorating and is solved as a one-dimensional root, in contrast to the multi-parameter fixed point a calibration to price moments would require.

\subsection{Analytical properties}\label{sub:model_propositions}

\begin{proposition}[Threshold]\label{prop:existence}
For $\sigma>0$, $0<R<1$, $\psi>0$, equation~\eqref{eq:F} has a unique root $s^* \in [\,1-\sigma,\,1+\sigma\,]$. Comparative statics follow from the implicit function theorem, $\partial s^*/\partial R = -F_R/F_{s^*}$.
\end{proposition}

\begin{proposition}[Endogenous stress probability]\label{prop:ps}
The probability of the stress regime, the event $D > R$ that the coin trades below par, is
\begin{equation}\label{eq:ps}
p_s = \Pr(\theta < \theta_s) = \Phi\!\left(\frac{\theta_s - 1}{\sigma_\theta}\right),
\qquad \frac{\partial p_s}{\partial R} = \frac{2\sigma}{\sigma_\theta}\,\phi\!\left(\tfrac{\theta_s-1}{\sigma_\theta}\right)\!\cdot\!\frac{\partial \theta_s}{\partial R} < 0,
\end{equation}
so deeper primary-redemption access lowers the run probability. This is the run probability of \citet{klee_etal_2026_fragility}, here endogenous and specific to each issuer rather than a frequency shared across coins.
\end{proposition}

\begin{proposition}[Reserve comparative static]\label{prop:reserves}
The expected depeg loss $\ell_D = \mathbb{E}_\theta[D(1-p^{sec})]\cdot 10^4$ is closed form, and $\partial \ell_D/\partial R < 0$ in the interior: deeper liquid reserves lower the expected depeg loss. Under a uniform fundamental prior the integral is polynomial; under the normal prior it is a truncated-moment expression in $\Phi$ and $\phi$.
\end{proposition}

The sign reads opposite to the bank-run intuition, and the difference is worth stating. In \citet{goldstein_pauzner_2005_demand_deposits} stronger fundamentals shrink the run region because withdrawal is destructive. Redemption here is conversion at a price, not destructive withdrawal: more par capacity makes exit safer, so the conversion mass weakly rises, but what deeper reserves reduce is the \emph{cost} of conversion, because a larger share clears at par and the spillover clears against a higher secondary price. The priced object is the loss, not the mass. A supervisor reads Proposition~\ref{prop:reserves} accordingly: liquidity requirements lower the holder's expected loss precisely by making redemption safe enough that more of it happens at par.

\subsection{The expected loss}\label{sub:what_it_delivers}

The model delivers, for each issuer, the stress depeg loss in basis points,
\begin{equation}\label{eq:ellDs}
\ell_{D,s} = \mathbb{E}_\theta\!\left[D(\theta; s^*)\,\bigl(1 - p^{sec}(D(\theta; s^*))\bigr)\right] \cdot 10^4,
\end{equation}
the probability-weighted depeg over the fundamental prior. The counterparty term $\ell_C = h\,\xi\cdot 10^4$ and the baseline depeg $\ell_{D,b}$, the cost of converting under normal conditions, are observable (Section~\ref{sec:risk_factors}): the linear model has zero baseline mid-depeg by construction, so $\ell_{D,b}$ is the observed off-ramp basis rather than a model output. The holding leg's expected loss is the sum
\begin{equation}\label{eq:total_risk_premium}
\Pi_{\text{risk}} = \underbrace{\ell_C}_{\text{observable}} + \underbrace{\ell_{D,b}}_{\text{observable basis}} + \underbrace{\ell_{D,s}}_{\text{closed-form game}},
\end{equation}
which we add to the on-ramp, on-chain, and off-ramp fees of the circuit. Section~\ref{sec:calibration} parametrises the model by observables and contrasts the closed form with the reduced-form evidence.
